\section{A uniform bound for the top-form Hodge structure}
\label{sec:rank}

The diagonal estimate bounds the top-form Hodge rank from a uniform
positive lower bound for the Seshadri constant of a normalized
polarization.  Here we remove that additional hypothesis for varieties
whose smaller rational images are rationally connected.

\begin{theorem}\label{thm:rank}
For every integer $n\ge2$ there is a constant $R_n$ with the following
property.  Let $T$ be a projective $\Q$-factorial terminal complex
$n$-fold with $K_T\sim0$.  Suppose that a smooth projective resolution
$U$ has $h^1(U,\OO_U)=0$, and that every positive-dimensional rational
image of $T$ of dimension less than $n$ has a rationally connected smooth
projective resolution.  Then
\[
                         r(T)=\dim_{\Q}T(U)\le R_n.
\]
\end{theorem}

The proof proceeds by contradiction.  Large $r(T)$ produces covering
subvarieties of Kodaira dimension zero on which complete systems of the
restricted polarization have small orders.  The rational functions constant on these
members define fields that can be realized by contractions on small
terminal models.  We choose a contraction with maximal proper base
dimension and use a bounded comparison model for its base.  A suitable
big divisor then lets us compare two section counts: vertical jets
bound the number of independent restrictions to a fibre, while slopes
on the base bound the number of global sections these restrictions can
produce.  Their product is too small to account for the divisor's
volume.

Unless another dependence is displayed, constants depend only on the
dimension.  Thresholds for taking a sufficiently divisible
multiple of a divisor may depend on that divisor and its variety.  We
write $D_1\preceq D_2$ when $D_2-D_1$ is rationally linearly equivalent
to an effective rational divisor; such a comparison is restricted only
to subvarieties not contained in the chosen effective difference.

\subsection{Section orders and the covering-family inputs}

Let $P$ be a big, nef, semiample rational Cartier divisor on an integral
projective variety $V$.  At a smooth point $x$, set
\[
 \gamma(P;V,x)=
 \sup_{\substack{m>0\text{ sufficiently divisible}\\
                 0\ne s\in H^0(\widetilde V,m\mu^*P)}}
                         \frac{\ord_x(s)}m,
\]
where $\mu:\widetilde V\to V$ is a smooth projective resolution that is
an isomorphism near $x$.  Birational pushforward makes this independent
of the resolution.  We always use these complete section spaces, even
when $V$ occurs as a subvariety of another variety.  Write
$\gamma(P;V)$ for the very general value.  In dimension $d$, counting
sections and jet conditions gives
\begin{equation}\label{eq:rank:order-volume}
                         \gamma(P;V)\ge(P^d)^{1/d}.
\end{equation}
Indeed, the two leading counts are $P^d m^d/d!$ and $k^d/d!$ for
sections and conditions of order at least $k$, respectively.

We use three results from the scalar-order construction of
\citet[Section ``Scalar order and curve estimates'']{UPI}.  We state
exactly the parts needed below.  A \emph{marked covering family} consists
of an integral parameter space $S$, a dominant mark map $u:S\to V$,
and an incidence $\mathcal I\subset S\times V$ whose geometric generic
fibre $C_b$ is integral, positive-dimensional, and contains $u(b)$.
A generic movement from a generic mark chooses a generic parameter over
that mark and then a generic point of its member.  Each new choice is
generic over all preceding data.  Dominant extensions of parameter
spaces are allowed to specify geometric components and models.

\begin{externalinput}[Generic chains]\label{input:rank:chains}
A finite nonempty collection of marked covering families determines a
rational fibration $\xi:V\dashrightarrow M$ with geometrically integral
generic fibre.  If its generic fibre, called a leaf, has dimension $h$,
then $1\le h\le\dim V$, every generic member lies in a leaf, and at most
$h$ generic movements reach a point generic in the leaf over the
algebraic closure of the field of its starting mark.  Suppose that,
for each family, either
\[
 \gamma(P|_{C_b};C_b,u(b))\le B,
 \qquad\text{or}\qquad
 \dim C_b=1,\quad
 \frac{P\cdot C_b}{\mult_{u(b)}C_b}\le B.
\]
For a smooth model $F$ of the geometric generic leaf, with $P$ pulled
back, one then has
\begin{equation}\label{eq:rank:chain-bound}
                  \gamma(P|_F;F)\le hB,
                  \qquad (P|_F)^h\le(hB)^h.
\end{equation}
\end{externalinput}

This is \cite[Lemma~6.2]{UPI}.
The assertion about endpoints also identifies its function field:
$\C(M)$ consists exactly of the rational functions constant on the
geometric generic members.  In fact, enlarge the field of definition
to contain the function being tested.  Constancy along successive
movements then gives constancy at a generic endpoint of the leaf,
hence constancy on that geometric generic leaf.

\begin{externalinput}[Tracking and canonical degrees]\label{input:rank:tracking}
Let $Z$ be projective, $\Q$-factorial and klt of dimension $d\ge2$, and
let $N$ be big, nef and semiample.  If $d+1$ equally spaced positive
rational levels, with spacing $\Delta>0$, lie strictly between
$(N^d)^{1/d}$ and $\gamma(N;Z)$, there is a covering family of
subvarieties $W\subset Z$, of dimension $0<e<d$, such that a smooth
geometric generic model $W^*$ satisfies
\begin{align}
 N^eW&\le(2/\Delta)^{d-e}N^d,\label{eq:rank:tracking-degree}\\
 K_{W^*}(N|_{W^*})^{e-1}
 &\le\bigl(K_Z+(2d/\Delta)N\bigr)|_W(N|_W)^{e-1}.
 \label{eq:rank:tracking-canonical}
\end{align}
If a smooth model of $Z$ has nonnegative Kodaira dimension, so does
$W^*$.  The same inheritance statement holds for arbitrary covering
families.
\end{externalinput}

This is \cite[Lemma~6.3]{UPI}; its section
order argument uses the parameter differentiation of
\citet[Proposition~2.3]{EKL1995}.  For the inheritance statement, cut the
parameter space until incidence evaluation is generically finite and
dominant, resolve, and restrict the ramification formula to a geometric
generic parameter fibre.  The cuts can be chosen with independent
generic coefficients, so the construction still tests the original
geometric generic member.

\begin{externalinput}[Small volumes on Iitaka fibres]\label{input:rank:small-fibre}
Fix $e\ge1$ and $C'>0$.  Let $W_i$ be smooth projective $e$-folds with
$\kappa(W_i)\ge0$, and let $L_i$ be big, nef, semiample rational
divisors such that
\[
                 L_i^e\longrightarrow0,
                 \qquad K_{W_i}L_i^{e-1}\le C'L_i^e.
\]
After passing to a subsequence and resolving their Iitaka fibrations,
the smooth geometric generic fibres $U_i$ have a fixed positive
dimension, satisfy $\kappa(U_i)=0$, and obey
\[
                    \vol(K_{U_i}+L_i|_{U_i})\longrightarrow0.
\]
The restrictions $L_i|_{U_i}$ remain big, nef and semiample.  When the
Iitaka base is a point, $U_i$ is the entire smooth model.
\end{externalinput}

This is \cite[Lemma~6.4]{UPI}.  Its uniform Iitaka hypothesis is supplied by that paper's
uniform pluricanonical Iitaka theorem.  We use the conclusion in
arbitrary lower dimensions, including a zero-dimensional Iitaka base.
All three inputs hold at geometric generic data: their finite
constructions spread after dominant parameter extensions.  A countable
field of definition includes the section-order conditions in all
sufficiently divisible degrees.  Thus the analytic and model-theoretic
inputs over $\C$ apply to these data without specializing a previously
generic parameter.

\subsection{Large rank produces subvarieties of small order}

The transition is indirect: small Seshadri constants force large
ambient section orders, and the tracking input turns these orders into
subvarieties of small polarized degree.  Passing to Iitaka fibres and repeating in
decreasing dimension eventually produces small complete-system orders
on Kodaira-dimension-zero members.

\begin{lemma}\label{lem:rank:low-order}
Fix $n\ge2$ and $\eta>0$.  There is a number $R(n,\eta)$ such that if
$T$ is as in Theorem~\ref{thm:rank} and $r(T)>R(n,\eta)$, the following
holds.  For every normal projective canonical $V$ birational to $T$
with $K_V\sim0$, and every ample rational Cartier divisor $P$ on $V$,
there is a marked covering family of subvarieties $C_b$ with
$0<\dim C_b<n$, whose smooth geometric generic models satisfy
\begin{equation}\label{eq:rank:low-order}
 \kappa(C_b^*)=0,
 \qquad
 \gamma(P|_{C_b};C_b,u(b))\le\eta(P^n)^{1/n},
\end{equation}
where the mark is smooth on the member.
\end{lemma}

\begin{proof}
We explain the two scalar-order arguments from \cite{UPI} used here,
since the additional fibration hypothesis in its full scalar-estimate
Proposition~6.1 is unnecessary for this conclusion.

\emph{From small Seshadri constants to large ambient orders.}
Suppose that $1\le P_i^n\le2$ and that the very general Seshadri
constants $\varepsilon(P_i)$ tend to zero.  Then
\begin{equation}\label{eq:rank:large-order}
                         \gamma(P_i;V_i)\longrightarrow\infty
\end{equation}
after passing to a subsequence.  Otherwise the orders are bounded by
some $D_0$.  Curves realizing arbitrarily small Seshadri ratios spread
into marked covering curve families.  Input~\ref{input:rank:chains}
produces leaves of a fixed dimension $h$ with $(P_i|_{F_i})^h\to0$.
Since $P_i^n\ge1$, eventually $0<h<n$.  Resolve the leaf fibration and
choose a very general smooth fibre $F_i$; put $q=n-h$.  Its normal
bundle is trivial.  Therefore restriction to its $u$th infinitesimal
neighbourhood has rank at most
\[
                 \binom{u+q-1}{q}h^0(F_i,mP_i|_{F_i}).
\]
Set $u=\lceil Dm\rceil$ for a fixed $D>D_0$.  For each fixed
fibration the leading coefficient of this expression in $m^n$ is
\[
                   \frac{D^q}{q!h!}(P_i|_{F_i})^h,
\]
whereas the ambient section count has coefficient at least $1/n!$.
First take $i$ large, then $m$ sufficiently divisible and large.  A
nonzero ambient section vanishes on that neighbourhood, so has order
at least $u$ at a very general point of $F_i$.  This contradicts the
bound $D_0$.  This argument is the proof of the second scalar inequality
in \cite{UPI}, with its order bound as an explicit hypothesis.

\emph{From large ambient orders to small orders on moving members.}
If the lemma were false, choose counterexamples with $r(T_i)\to\infty$
and rescale $P_i$ rationally to have $1\le P_i^n\le2$.  Lemma
\ref{lem:diagonal} forces $\varepsilon(P_i)\to0$, and hence
\eqref{eq:rank:large-order} holds.  Take small projective
$\Q$-factorializations and pull back $P_i$.  They are klt, have
principal canonical divisor, and preserve the complete section spaces.
Choose $n+1$ order levels with spacing $\Delta_i\to\infty$ between
$(P_i^n)^{1/n}$ and $\gamma(P_i;V_i)$.  Input
\ref{input:rank:tracking} produces members of a fixed dimension $e<n$
on a subsequence.  With $W_i$ their smooth models and $L_i=P_i|_{W_i}$,
\[
 \kappa(W_i)\ge0,\qquad
 L_i^e\longrightarrow0,\qquad
 K_{W_i}L_i^{e-1}\le C'L_i^e.
\]
Here $K_{V_i}\sim0$, so $C'$ can be fixed.  Input
\ref{input:rank:small-fibre} gives positive-dimensional Iitaka fibres
$U_i$ with $\kappa(U_i)=0$ and adjoint volume tending to zero.

Represent $L_i|_{U_i}$ by a divided general free member
$\Gamma_i$ for which $(U_i,\Gamma_i)$ is klt.  A $\Q$-factorial good
model $Z_i'$ of this big adjoint exists by \cite{BCHM2010}; write $N_i'$
for its big semiample adjoint.  Then
\[
 (N_i')^{\dim Z_i'}\longrightarrow0,
 \qquad K_{Z_i'}\preceq N_i',
 \qquad\kappa(Z_i')=0.
\]
On a common resolution, with maps $x_i$ to $U_i$ and $y_i$ to $Z_i'$,
there is also the effective comparison
\begin{equation}\label{eq:rank:polarization-comparison}
                         x_i^*(L_i|_{U_i})\preceq y_i^*N_i'.
\end{equation}
To prove it, write
$x_i^*(K_{U_i}+\Gamma_i)\sim_{\Q}y_i^*N_i'+F_i$ with $F_i\ge0$
exceptional over $Z_i'$.  Since $\kappa(U_i)=0$, choose an effective
$D_i\sim_{\Q}K_{U_i}$.  In divisible degrees, the pullback of $D_i$
plus a general member of the free polarization system contains the
fixed part $F_i$.  The latter general member contains none of the
finitely many components of $F_i$, hence $F_i\le x_i^*D_i$.  Subtracting
proves \eqref{eq:rank:polarization-comparison}.

If $\gamma(N_i';Z_i')\to0$ along a subsequence, stop.  Otherwise take
a subsequence on which it has a positive lower bound.  More generally,
the data at this stage have fixed dimension $d$, are $\Q$-factorial klt,
and satisfy
\[
 N_i^d\to0,\qquad \kappa(Z_i)\ge0,\qquad
 K_{Z_i}\preceq a_0N_i,\qquad \gamma(N_i;Z_i)\ge a_1>0
\]
for fixed positive $a_0,a_1$.  Choose $d+1$ equally spaced fixed levels
below $a_1$; they eventually exceed $(N_i^d)^{1/d}$.  The tracking and
Iitaka-fibre inputs apply again, since their canonical-degree bound is
at most $(a_0+2d/\Delta)L_i^e$.  They produce another good model of
strictly smaller positive dimension, with polarization volume tending
to zero, Kodaira dimension zero, canonical class bounded above by the
new polarization, and comparison
\eqref{eq:rank:polarization-comparison}.  At this new stage we make the
same dichotomy: either the orders tend to zero along a subsequence and
we stop, or they have a positive lower bound on a subsequence and the
construction repeats.  Every repetition decreases dimension.  In
dimension one the order is bounded by the degree, which tends to zero,
so the construction must terminate with order tending to zero.

Trace these members through the successive covering families and
Iitaka fibres to the original $V_i$.  Fibres sweep the preceding member;
the birational comparisons are isomorphisms on opens met by the next
geometric generic members.  Thus the composite members map birationally
to their images and have Kodaira dimension zero.  Each effective
comparison restricts to the next member, which is not contained in its
extra divisor.  Restricting successively and pulling back to a common
smooth model of the final descendant compares the original restricted
polarization with the final polarization by an effective difference.
Multiplication by that difference injects their \emph{complete} section
spaces.  At a generic mark outside the difference it preserves order.
Thus the resulting members satisfy
$\gamma(P_i|_{C_b};C_b,u(b))\to0$.  They eventually satisfy
\eqref{eq:rank:low-order}, a contradiction.  This is the
dimension-decreasing part of the first scalar-estimate proof in
\cite{UPI}, stopped before its separate fibration contradiction.
\end{proof}

We will also apply this lemma to a big rational divisor $D$ on a
$\Q$-factorial canonical model $X$ birational to $T$, with $K_X\sim0$.
Choose a small effective klt boundary proportional to $D$ and run its
MMP.  Its ample model $X_D$ is canonical with $K_{X_D}\sim0$: the
canonical generator remains principal on every model, and all these
birational steps are crepant for the zero boundary.  Let $P_D$ be the
ample polarization induced by $D$.  On a common resolution,
\begin{equation}\label{eq:rank:ample-model}
 p^*D\sim_{\Q}q^*P_D+E,
 \qquad E\ge0\text{ exceptional over }X_D,
 \qquad v(D):=\vol(D)^{1/n}=(P_D^n)^{1/n}.
\end{equation}
Multiplication by the fixed part identifies sufficiently divisible
section spaces.  The low-order lemma applies to $(X_D,P_D)$ because
$r(T)$ is birationally invariant.

\subsection{Realizing chain fields by contractions}

The chain fibration just constructed is initially rational.  Its
realization as a morphism on a small terminal model will make its base
accessible to the boundedness input.  Call a subfield of $\C(T)$
\emph{attained} if it is the field of a contraction $T'\to Y$, where
$T'$ is $\Q$-factorial terminal and is equipped with a crepant birational
map to $T$ that is an isomorphism in codimension one.

\begin{lemma}\label{lem:rank:attainment}
If $0<\eta<1/n$, every chain field obtained from a family in
Lemma~\ref{lem:rank:low-order} is attained and has transcendence degree
strictly between $0$ and $n$.
\end{lemma}

\begin{proof}
Its leaves have positive dimension.  A leaf cannot fill $V$, since
\eqref{eq:rank:chain-bound} and \eqref{eq:rank:low-order} would give
$\gamma(P;V)\le n\eta(P^n)^{1/n}$, contradicting
\eqref{eq:rank:order-volume}.  Its function field is relatively
algebraically closed in $\C(T)$ by the geometric integrality of its
leaves.

Choose a smooth projective base model $M$ for this field.  For
$\phi:T\dashrightarrow M$, let $J$ be the closure on $T$ of the inverse
image of a general member $H_M$ of a very ample system on $M$.
The map is defined in codimension one, and these divisors form a movable
system.  On a smooth resolution $p:\widetilde T\to T$ resolving $\phi$,
\begin{equation}\label{eq:rank:movable-base}
            p^*J=\widetilde\phi^*H_M+E_M,
            \qquad E_M\ge0\text{ exceptional over }T.
\end{equation}
Choose $H_M$ general for this resolution as well.  With the canonical
generator used to represent $K_{\widetilde T}$, terminality gives
\begin{equation}\label{eq:rank:terminal-domination}
                         K_{\widetilde T}\ge aE_M
\end{equation}
for some positive rational $a$.

We claim that all ratios of sections of divisible multiples of $J$ are
constant on the generic members of the covering family, transferred to
$\widetilde T$.  Here is why \eqref{eq:rank:terminal-domination} may be
used on a member.  Cut the parameter space by the dimension of a generic
evaluation fibre, choosing hyperplanes with independent generic
coefficients.  The resulting incidence evaluates generically finitely
and dominantly to $\widetilde T$.  Nonempty zero-dimensional cuts of the
generic evaluation fibre avoid the added boundary in a projective
closure.  The universal choices of hyperplanes still dominate the
original parameter space.  After a geometric component and a resolution
are chosen, the ramification formula on this total space, restricted to
its smooth geometric generic parameter fibre $C^*$, gives
\[
                             K_{C^*}\succeq aE_M|_{C^*}.
\]
The map $\widetilde\phi$ is constant on $C^*$, so
\eqref{eq:rank:movable-base} identifies the restricted polarization with
$E_M|_{C^*}$.  Two independent restricted sections in any common degree,
followed by a suitable power and multiplication by the effective
canonical difference, would give two independent sections of a
pluricanonical system on $C^*$.  That would contradict $\kappa(C^*)=0$.
Ambient nonzero sections remain nonzero at these generic data.  The
claim follows.

By the function-field characterization after
Input~\ref{input:rank:chains}, the section-ratio field of $J$ is contained
in $\C(M)$.  The reverse inclusion follows from ratios of hyperplane
pullbacks, since $H_M$ is very ample.  Thus this section-ratio field is
exactly $\C(M)$.

For a small positive rational $\varepsilon$, the pair
$(T,\varepsilon J)$ is klt and its adjoint is pseudo-effective.  Theorem~\ref{input:good-model} gives a terminating MMP to a good model.  There
are no divisorial steps: the divisor is movable, whereas the exceptional
divisor of a negative divisorial contraction is contained in its fixed
part.  Movability persists through the flips.  These steps are crepant
for the zero boundary; the resulting $T'$ is again terminal and
isomorphic to $T$ in codimension one.  The semiample contraction, with
Stein factorization, realizes the section-ratio field $\C(M)$, which is
already relatively algebraically closed.  This proves attainment.
\end{proof}

Assume now that $r(T)>R(n,\eta)$, where $\eta$ will be chosen uniformly
at the end.  There is an attained field of transcendence degree between
$0$ and $n$.  Among all attained fields of transcendence degree less
than $n$, choose one whose transcendence degree $b$ is maximal, and
write its contraction as
\[
                    h:T'\longrightarrow Y,
                    \qquad s=n-b>0.
\]
The rational-image hypothesis in Theorem~\ref{thm:rank} makes a smooth
resolution of $Y$ rationally connected.  Proposition~\ref{input:bounded-base}
therefore provides a normal comparison model $\widehat Y$, isomorphic
to $Y$ in codimension one, and a very ample Cartier divisor $H$ with
\begin{equation}\label{eq:rank:base-degree}
          N:=H^b\le C(n),
          \qquad \abs{K_{\widehat Y}H^{b-1}}\le C(n).
\end{equation}

We make one further modification so that divisors pulled back from the
common big opens of these bases extend without ambiguity upstairs.
Let $F_0$ be the sum of the prime divisors of $T'$ whose image in $Y$
has codimension at least two.  There are finitely many: they are among
the components of the closed locus where the fibre dimension is larger
than $s$.  Choose a very ample $H_0$ on $Y$ and a rational number
$a>2n$.  Run a terminating MMP for a klt boundary rationally equivalent
to
\[
                        \varepsilon F_0+a h^*H_0,
\]
with $\varepsilon>0$ small; the latter summand is represented by divided
general members.  The length bound $2n$ for negative extremal rays,
applied to the pair omitting this summand, forces every step to be over
$Y$.  Indeed a horizontal extremal curve has $h^*H_0$-degree at least
one, so the added summand makes its total degree positive.  This
argument repeats after each step, and the morphism to $Y$ descends
through contractions and flips.

Every component of $F_0$ is contracted.  To see this, over the complement
of their images the sections of the displayed adjoint come from $aH_0$:
the fibres are connected and the adjoint restricts trivially there.
These base sections extend across the omitted codimension-two set by
normality.  Thus every surviving component of $F_0$ would be a fixed
component in all sufficiently divisible systems, contradicting
semiampleness on the good model.  Write that model as
\begin{equation}\label{eq:rank:working-fibration}
                              h:X\longrightarrow Y.
\end{equation}
It is $\Q$-factorial canonical, has $K_X\sim0$, and defines the same
field as before.  Its resolution still has irregularity zero.  Crucially,
\emph{no prime divisor of $X$ maps to a subset of codimension at least
two in $Y$}.  Terminality is no longer required for this working model;
any later chain field is attained by applying
Lemma~\ref{lem:rank:attainment} on the original terminal $T$.

\subsection{A big divisor calibrated by the degree on the base}
\label{subsec:rank:calibration}

To compare orders on fibres with slopes on the base, we need a linear
functional on divisor classes of $X$ that agrees with base degree on
pullbacks and is nonnegative on effective divisors.  We construct this
functional together with a big divisor whose volume root and functional
value are controlled by one parameter.  The base polarization alone has
zero volume on $X$, whereas an arbitrary ample divisor need not satisfy
the required degree bound.

\begin{lemma}[Calibration]\label{lem:rank:calibration}
Let $h:X\to Y$ be a contraction of normal projective complex varieties,
where $X$ is a $\Q$-factorial canonical $n$-fold with $K_X\sim0$ and
irregularity zero on a resolution.  Suppose that $0<b:=\dim Y<n$ and
that no prime divisor of $X$ maps into a subset of codimension at least
two in $Y$.  Let $\widehat Y$ be a normal projective variety isomorphic
to $Y$ in codimension one, and let $H$ be a very ample divisor on
$\widehat Y$ satisfying
\[
 H^b\le N_0,
 \qquad
 \bigl|K_{\widehat Y}\cdot H^{b-1}\bigr|\le N_0.
\]
Pullback from the common smooth open subsets of the bases defines
\[
 W=h^*\Cl(Y)_{\R}
   =h^*\Cl(\widehat Y)_{\R}
   \subset N^1(X)_{\R},
 \qquad
 \ell(h^*L)=L\cdot H^{b-1}.
\]
Let $A=h^*H$, represented by an effective integral Weil divisor.
There is a linear extension
$\ell_X:N^1(X)_{\R}\to\R$ of $\ell$, nonnegative on the
pseudo-effective cone, and there are a rational number $\delta>0$ and
a big rational divisor $D$ such that
\begin{equation}\label{eq:rank:calibration}
 c\delta\le v(D)\le C\delta,
 \qquad
 D-\frac{\delta}{3}A\text{ is big},
 \qquad
 \ell_X(D)\le C\delta.
\end{equation}
Here $v(D)=\vol(D)^{1/n}$, $c>0$ depends only on $n$, and $C$ depends
only on $n,N_0$.
\end{lemma}

The two conditions beyond $v(D)\asymp\delta$ have distinct roles.
The bigness of $D-\delta A/3$ will force the low-order families to be
vertical.  The bound on $\ell_X(D)$ will convert cotangent
semipositivity into a slope bound on the base.

\begin{proof}
We first construct a supporting functional for the pseudo-effective
cone and then choose a point on which volume has the required scale.
The hypothesis on irregularity, together with $\Q$-factoriality, gives
\[
 \Cl(X)_{\R}=\Pic(X)_{\R}=N^1(X)_{\R}.
\]
Indeed pullback to a resolution detects numerically trivial line
bundles, whose classes are torsion when the Picard variety is zero.
We can therefore pass between real linear equivalence and numerical
classes throughout the proof.

Choose common smooth big open subsets of $Y$ and $\widehat Y$.
Their inverse image in $X$ has complement of codimension at least two,
so Cartier pullbacks there extend uniquely as Weil divisors on $X$.
This defines the indicated pullback on class groups.  It is injective:
if $s=n-b$ and $Q$ is an ample Cartier divisor on $X$, the operation
\[
 B\longmapsto h_*(B\cdot Q^s)
\]
on divisor classes is a left inverse to pullback, up to multiplication
by the positive degree of $Q^s$ on the generic fibre.  The same
operation sends effective classes to effective classes or zero, by
using general representatives of sufficiently divisible multiples of
$Q$.

The degree $\ell$ is strictly positive on
$W\cap\overline{\operatorname{Eff}}(X)\setminus\{0\}$.
To see the point requiring care, choose a finite linear projection
$f:\widehat Y\to\PP^b$ with $f^*\OO_{\PP^b}(1)=\OO_{\widehat Y}(H)$.
For every effective real Weil divisor $L$ on $\widehat Y$,
\begin{equation}\label{eq:rank:degree-dominates}
 \ell(L)H-L\quad\text{has an effective representative}.
\end{equation}
In fact $f^*f_*L-L\ge0$ and
$f_*L\sim_{\R}\ell(L)\OO_{\PP^b}(1)$.
The intersect-and-push operation just described shows that if a
pullback class is pseudo-effective on $X$, its base class lies in the
closure of effective base classes.  Applying
\eqref{eq:rank:degree-dominates} to approximating classes and pulling
back shows that a degree-zero such class and its negative are both
pseudo-effective on $X$.  The pseudo-effective cone in $N^1(X)_{\R}$
contains no line, so the class is zero.  This also proves that the
formula defining $\ell$ is well defined on $W$.

Fix an ample class $P_0$ on $X$, and minimize $\ell(L)$ on the set
\[
 \{L\in W:P_0+L\in\overline{\operatorname{Eff}}(X)\}.
\]
The set is nonempty.  Its subsets on which $\ell$ is bounded above
are bounded: otherwise, dividing an unbounded sequence by its norm
and passing to a limit would give a nonzero pseudo-effective class in
$W$ of nonpositive degree.  Closedness therefore gives a minimum
$c_0=\ell(L_0)$.  Put $d_0=P_0+L_0$.
The open big cone is disjoint from
\[
 P_0+\{L\in W:\ell(L)\le c_0\}.
\]
Indeed a big point in this affine half-space would remain big after
subtracting a sufficiently small positive multiple of $A$, contradicting
minimality.  Separating these two convex sets gives a nonzero linear
functional $\lambda$, nonnegative on the pseudo-effective cone, with
$\lambda(d_0)=0$.  Its restriction to $W$ is a nonnegative multiple
of $\ell$: it vanishes on $\ker\ell$, and the orientation follows
from the half-space being separated.  That multiple is positive,
since otherwise $\lambda(P_0)=\lambda(d_0)=0$, whereas a nonzero
nonnegative functional is strictly positive at an interior point of
the cone.  Normalize $\lambda$ to obtain $\ell_X$.

For $t>0$, set $d_t=d_0+tA$.  These classes are big.  For large $t$,
$L_0+tA$ is pseudo-effective, since sufficiently large multiples of
$H$ added to any fixed real Weil class on the base are effective;
convexity with the pseudo-effective class $d_0$ then proves the
assertion for every $t>0$.  On the other hand, $A$ is not big: its
restriction to the positive-dimensional generic fibre is trivial.
Continuity and homogeneity of volume give
\begin{equation}\label{eq:rank:large-calibration-parameter}
 \frac{v(d_t)}t=v(A+t^{-1}d_0)\longrightarrow0
 \qquad(t\longrightarrow\infty),
 \qquad
 \ell_X(d_t)=tH^b.
\end{equation}
It remains to bound this volume ratio from below for small $t$.
This is where the trivial canonical class enters the construction.

Let $p:\widetilde X\to X$ be a smooth projective resolution and write
the divisorial Zariski decomposition of
$\widetilde d_0=p^*d_0$ as
\[
 \widetilde d_0=P_\sigma(\widetilde d_0)
                      +N_\sigma(\widetilde d_0).
\]
We use its numerical formulation, as in
\citep[Chapter~III]{Nakayama2004} and \citep{Boucksom2004}.
Fix the dimension-dependent integer $M=M(n)$ in Birkar's
Calabi--Yau effective birationality theorem
\citep[Corollary~1.4]{Birkar2023}.
Suppose first that $P_\sigma(\widetilde d_0)$ is numerically nonzero.
For an ample class $J$ on $\widetilde X$, monotonicity of positive
intersection products gives
\begin{align*}
 \vol(\widetilde d_0+uJ)
 &=\bigl\langle(\widetilde d_0+uJ)^n\bigr\rangle\\
 &\ge u^{n-1}
 P_\sigma(\widetilde d_0+uJ)\cdot J^{n-1}.
\end{align*}
Here the trace of the one-class positive product is the positive
part of the divisorial Zariski decomposition; see
\citep[Definition~2.10, Proposition~2.13, Theorem~3.1, and
Section~3.4]{BFJ2009}.
As $u\downarrow0$, the intersection on the right tends to
$P_\sigma(\widetilde d_0)\cdot J^{n-1}>0$.
Choose $a>0$ such that $p^*d_1-aJ$ is pseudo-effective.  For
$0<t<1$, the identity
$d_t=(1-t)d_0+td_1$ and volume monotonicity yield
\[
 \vol(d_t)\ge(1-t)^n
 \vol\left(\widetilde d_0+\frac{at}{1-t}J\right).
\]
Consequently $v(d_t)/t\to\infty$ as $t\downarrow0$.
The constants in this divergence may depend on $X$; no uniform
threshold in $t$ is needed.

Suppose instead that $P_\sigma(\widetilde d_0)$ is numerically zero.
Then $d_0$ has a unique effective real representative $N_0'$.
For existence, push forward $N_\sigma(\widetilde d_0)$; numerical
and real linear equivalence agree here.  For uniqueness, the pullback
of any effective representative of $d_0$ dominates
$N_\sigma(\widetilde d_0)$.  Their difference is effective and
numerically zero, hence zero.  Let $E_0$ be the reduced support of
$N_0'$.  It has Kodaira dimension zero: if some multiple of $E_0$ had
a different effective representative, replacing a sufficiently small
positive multiple of $E_0$ inside $N_0'$ would contradict uniqueness.

Choose a small positive rational $\varepsilon$ for which
$(X,\varepsilon E_0)$ is klt.  Theorem~\ref{input:good-model}
gives a terminating MMP for this pair and a
good model $\overline X$.  Precisely the components of $E_0$ are
contracted.  Any divisorial step must contract a component of its
transform, since $K\sim0$ and a curve outside the effective boundary
has nonnegative intersection with it.  Conversely, semiampleness on
the good model and $\kappa(E_0)=0$ force the effective pushforward
of $E_0$ to vanish.  These steps are crepant for the zero boundary, so
$\overline X$ is again $\Q$-factorial canonical with principal
canonical divisor.  The pushforward $\overline A$ of $A$ is an
integral Weil divisor and is big: the big class $d_t$ pushes forward
to $t\overline A$.  Birkar's theorem now gives
\begin{equation}\label{eq:rank:integral-volume-lower-bound}
 \vol(\overline A)\ge M^{-n}
 \qquad\text{\citep[Corollary~1.4]{Birkar2023}}.
\end{equation}
Its hypotheses are exactly that $\overline X$ is projective klt,
$K_{\overline X}\sim_{\R}0$, and $\overline A$ is big and integral;
no Cartier-index bound is required.

A single rational number $u\ge0$ gives inclusions
\[
 H^0(\overline X,m\overline A)
 \subseteq H^0(X,m(A+uE_0))
\]
for every sufficiently divisible $m$.  Indeed, on a common resolution
pull back $\overline A$ and push forward to $X$.  The resulting
rational divisor differs from $A$ only on $E_0$, and its difference
is bounded above by $uE_0$.  Thus
$\vol(A+uE_0)\ge\vol(\overline A)$.
For all sufficiently small $t>0$ we have $N_0'\ge tuE_0$, and hence
\[
 \vol(d_t)=\vol(N_0'+tA)
 \ge t^n\vol(A+uE_0)\ge t^n M^{-n}.
\]
The same conclusion includes $E_0=0$, with $u=0$.

In both cases the ratio $v(d_t)/t$ is at least $M^{-1}$ for all
sufficiently small $t>0$.  By
\eqref{eq:rank:large-calibration-parameter} and continuity, it takes
the value $(2M)^{-1}$ at some $t>0$.  At this parameter,
\[
 d_t-\frac t3A=d_{2t/3}\text{ is big},
 \qquad \ell_X(d_t)=tH^b\le N_0t.
\]
Approximate $t$ by a positive rational number $\delta$ and $d_t$ by a
rational class $D$.  Openness of the big cone and continuity of volume
and $\ell_X$ preserve these assertions with, for example, lower
constant $1/(4M)$ and suitable upper constant depending only on
$n,N_0$.  Since $X$ is $\Q$-factorial, the rational class has a
rational Cartier divisor representative.  This proves
\eqref{eq:rank:calibration}.
\end{proof}

\subsection{Vertical orders and horizontal section counts}
\label{subsec:rank:count}

We now apply the low-order families to the divisor supplied by
Lemma~\ref{lem:rank:calibration}.  Keep the contraction $h:X\to Y$, its
bounded comparison model $\widehat Y$, and the very ample divisor $H$ on
$\widehat Y$.  Write $b=\dim Y$, $s=n-b$, $A=h^*H$, and $N=H^b$.
Thus $0<b<n$, and both $N$ and
$|K_{\widehat Y}\cdot H^{b-1}|$ are bounded in terms of $n$.
There is no divisor of $X$ mapping into the codimension-two sets omitted
when identifying $Y$ with $\widehat Y$.
Let $D$ and $\delta>0$ be the calibrated divisor and parameter.  In this
subsection $C$ denotes constants depending only on $n$, independently of
$0<\eta\le1$; the threshold for a sufficiently large divisible integer
$m$ may depend on all the fixed geometric data.

For sufficiently large divisible $m$, our target is the estimate
\begin{equation}\label{eq:rank:count-target}
 h^0(X,\OO_X(mD))
 \le C\bigl(m\eta v(D)\bigr)^s(m\delta)^b.
\end{equation}
The first factor will bound, over the base function field, the number
of independent restrictions to the generic fibre, using vertical jets.
The second will bound their
contribution to global sections, using slopes on the base and a section
estimate linear in the sheaf rank.  Since $v(D)$ is comparable to
$\delta$, this would contradict the volume asymptotic
$h^0(X,mD)\sim m^n v(D)^n/n!$ when $\eta$ is sufficiently small.

\paragraph{The families are vertical.}
Apply Lemma~\ref{lem:rank:low-order} to the ample model $X_D$ of $D$.
If a geometric generic member $C_t$ has nonconstant image in
$\widehat Y$, choose a hyperplane in $|H|$ through the image of its marked
point $x_t$ without containing the image of $C_t$.  Since
$D-\delta A/3$ is big, choose a section of a sufficiently divisible
multiple of this divisor whose divisor avoids the generic mark.  This
section,
multiplied by the corresponding power of the hyperplane section, gives
\[
  \gamma(P_D|_{C_t};C_t,x_t)\ge \delta/3.
\]
Here one works on a common resolution and uses
\eqref{eq:rank:ample-model}.  The mark lies in the birational isomorphism
locus and avoids the fixed divisor; the chosen section remains nonzero
on the member.  Hyperplane pullbacks extend in codimension one on $X$
because no divisor maps into the omitted base sets.  The same
construction is valid over the geometric generic coefficient field.
The low-order bound is at most $\eta v(D)\le C\eta\delta$, contradicting
this inequality when $\eta$ is sufficiently small in terms of $n$.

The family is therefore vertical for $h$.  Its chain field contains
$\C(Y)$ and is attained.  Both fields are relatively algebraically
closed in $\C(X)$; maximality of $\dim Y$ among proper attained fields
forces them to coincide.  Indeed the chain field is proper by
Input~\ref{input:rank:chains} and the choice $\eta<1/n$, and an inclusion
of equal transcendence degree would be algebraic.  On a smooth model
$F$ of the geometric generic fibre of $h$, also resolving the ample-model
map, the chain estimate \eqref{eq:rank:chain-bound} now gives
\begin{equation}\label{eq:rank:vertical-order}
  \gamma(P_D|_F;F)\le s\eta v(D).
\end{equation}

For sufficiently divisible $m$, put $V_m=H^0(X,\OO_X(mD))$ and let $r_m$
be the dimension over $\C(Y)$ of its span in the generic fibre of
$h_*\OO_X(mD)$.  After extending to the geometric generic field, this
span identifies with a subspace of $H^0(F,mP_D|_F)$ multiplied by the
fixed part in \eqref{eq:rank:ample-model}.  At a general mark that fixed
part is a unit.  Thus \eqref{eq:rank:vertical-order} bounds the order of
every nonzero element of the span, including arbitrary scalar
combinations over that field.  Its vertical $k$-jet map is injective for
\[
  k=\lceil ms\eta v(D)\rceil.
\]
Consequently, for $m$ sufficiently large,
\begin{equation}\label{eq:rank:vertical-rank}
  r_m\le {s+k\choose s}\le C\bigl(m\eta v(D)\bigr)^s.
\end{equation}
Saturate the subsheaf generically spanned by $V_m$ in $h_*\OO_X(mD)$,
transfer it across the common big open of $Y$ and $\widehat Y$, and extend
reflexively.  This gives a reflexive sheaf $\mathcal E_m$ on
$\widehat Y$ of rank $r_m$, with
\begin{equation}\label{eq:rank:sections-in-span}
  \dim V_m\le h^0(\widehat Y,\mathcal E_m).
\end{equation}
We must bound the number of these sections without losing a further
power of $r_m$.

\paragraph{A slope bound from jets along the fibres.}
For a torsion-free sheaf $\mathcal S$ on $\widehat Y$ write
\[
  \mu_H(\mathcal S)
  =\frac{c_1(\mathcal S)\cdot H^{b-1}}{\rank\mathcal S}.
\]
Let $\mathcal S\subset\mathcal E_m$ be saturated of rank $d>0$.
On a smooth big open of $X$, the saturated relative tangent sheaf
$\mathcal T=\ker(dh)$ and its quotient in $\mathcal T_X$ are vector
bundles.  We remove their codimension-two defects, while retaining the
generic points of divisors where the differential of $h$ drops rank.
The evaluation map from $h^*\mathcal S$ to
$L=\OO_X(mD)$ is defined there.  Since $\mathcal T$ is integrable, the
bundles of jets along $\mathcal T$ have a filtration with quotients
\[
  L\otimes\Sym^i(\mathcal T^*),\qquad 0\le i\le k.
\]
For completeness, these bundles can be constructed by taking ordered
derivatives in local frames of $\mathcal T$.  The product rule makes
changes of frames triangular in the order; integrability identifies the
degree-$i$ symbol with $\Sym^i(\mathcal T^*)$.  The cocycle identities
agree with ordinary fibrewise jets on the dense smooth locus of $h$
and hence hold throughout this open.

Functions from the base have zero derivative along $\mathcal T$.
Taking jets of the evaluation map therefore gives an $\OO_X$-linear map
\[
  h^*\mathcal S\longrightarrow J^k_{\mathcal T}(L),
\]
generically injective by \eqref{eq:rank:vertical-order}.  Take its top
exterior power and a nonzero component in the induced graded
filtration.  In characteristic zero, the symmetric and exterior powers
occurring there embed in direct sums of tensor powers.  Selecting a
nonzero summand gives
\begin{equation}\label{eq:rank:jet-determinant}
  h^*\det\mathcal S\longrightarrow
       (\mathcal T^*)^{\otimes j}\otimes L^d,
  \qquad 0\le j\le kd.
\end{equation}

Wedge with the pulled-back base differentials defines a regular map
\begin{equation}\label{eq:rank:base-wedge}
  \mathcal T^*\longrightarrow
  \bigwedge^{b+1}\Omega_X^1\otimes\OO_X(-h^*K_{\widehat Y}),
\end{equation}
which is generically injective.  To check regularity, extend a local
leafwise covector to an ordinary covector and wedge it with the pulled-back
base volume form.  Two extensions differ by a conormal covector, whose
wedge is zero.  This also proves regularity along multiple-fibre
divisors: the pulled-back volume form may vanish there but has no pole.

Combining \eqref{eq:rank:jet-determinant} and
\eqref{eq:rank:base-wedge}, and moving the twists to the source, gives a
nonzero map from a line bundle of class
\begin{equation}\label{eq:rank:source-line}
  h^*\det\mathcal S-dmD+jh^*K_{\widehat Y}
\end{equation}
to a tensor power of $\Omega_X^1$.  We test this map by the functional
$\ell_X$ from Lemma~\ref{lem:rank:calibration}.  On a smooth resolution
$p:Z\to X$, define $\Lambda(M)=\ell_X(p_*M)$.  This is nonnegative on
pseudo-effective divisor classes and vanishes on exceptional classes.
The canonical class of $Z$ is effective and exceptional, since $X$ is
canonical and $K_X\sim0$, so $\Lambda(K_Z)=0$.

Saturate the line defined by the map on $Z$.  Campana--P\u{a}un's
cotangent theorem \cite[Theorem~1.3]{CP2019} says that every torsion-free
quotient of a cotangent tensor power on $Z$ has pseudo-effective
determinant.  The determinant of the tensor power itself is a multiple
of $K_Z$.  It follows that the saturated line has $\Lambda$-degree at
most zero.  Its pushforward differs from \eqref{eq:rank:source-line}
by an effective divisor, because the original map is regular in
codimension one.  Thus
\begin{align}
  \ell(\det\mathcal S)
   &\le dm\ell_X(D)-j\ell(K_{\widehat Y})
        \notag\\
   &\le Cdm\delta.\label{eq:rank:slope-bound}
\end{align}
The last inequality uses $j\le kd$, $\eta\le1$, the bounded canonical
degree of $\widehat Y$, and $v(D)\le C\delta$; the ceiling in $k$ is
absorbed once $m\delta\ge1$.  If $j=0$, the same conclusion follows
directly from the effective divisor of the map to $\OO_X$.
Since $\mathcal S$ was arbitrary, this proves
$\mu_{\max,H}(\mathcal E_m)\le Cm\delta$.

\paragraph{Counting sections without a rank loss.}
The following elementary consequence of the Grauert--M\"ulich--Spindler theorem
removes the rank from inside the polynomial section bound.

\begin{lemma}\label{lem:rank:section-count}
Let $b\ge1$.  A torsion-free slope-semistable sheaf $\mathcal G$ on
$\PP^b_{\C}$, of rank $u$ and slope $\mu$ with respect to
$\OO_{\PP^b}(1)$, satisfies
\[
  h^0(\PP^b,\mathcal G)
    \le C_bu\bigl(1+\max\{0,\mu\}\bigr)^b,
\]
where $C_b$ depends only on $b$.
\end{lemma}

\begin{proof}
If $\mu<0$, a nonzero section would yield a subsheaf of slope at least
zero, contradicting semistability.  Suppose $\mu\ge0$.
For $b\ge2$, the Grauert--M\"ulich--Spindler theorem
\cite{GrauertMulich1975,Spindler1979}, in its torsion-free form
\cite[Corollary~3.1.6]{HuybrechtsLehn2010}, bounds the highest degree $a$
of the splitting on a general line by $a\le\mu+(u-1)/2\le\mu+u$.
For $b=1$, the splitting theorem on $\PP^1$ and
semistability give the stronger equality $a=\mu$.

A general linear flag gives, for any torsion-free sheaf with highest
general-line splitting degree $a$,
\begin{equation}\label{eq:rank:preliminary-count}
  h^0(\mathcal G)\le u{a+b\choose b}\quad(a\ge0),
  \qquad h^0(\mathcal G)=0\quad(a<0).
\end{equation}
Indeed the successive restrictions can be chosen torsion-free: embed
the sheaf in a vector bundle and choose the cuts general for the
associated primes of the quotient as well.  Testing general lines
shows that twists by integers below $-a$ have no sections.  Starting
with the bound $u(a+1)$ on a line, the hyperplane restriction sequences
and the identity
$\sum_{i=0}^a {i+b-1\choose b-1}={a+b\choose b}$ prove
\eqref{eq:rank:preliminary-count} inductively.

For an integer $t\ge1$, let
\[
  f_t:\PP^b\longrightarrow\PP^b,
  \qquad [x_0:\cdots:x_b]\longmapsto[x_0^t:\cdots:x_b^t].
\]
Finite separable pullback preserves semistability
\cite[Lemma~2.1]{Weissmann2024}.  Since $f_t$ is flat,
$\mathcal G_t=f_t^*\mathcal G\otimes\OO_{\PP^b}(b(t-1))$ is
torsion-free and semistable, of rank $u$ and slope
$t\mu+b(t-1)$.  Moreover
\begin{equation}\label{eq:rank:power-sections}
  h^0(\mathcal G_t)\ge t^b h^0(\mathcal G).
\end{equation}
To see this, multiply the pulled-back sections by the $t^b$ homogeneous
monomials
\[
  x_0^{\,b(t-1)-\sum_{i=1}^b e_i}
  \prod_{i=1}^b x_i^{e_i},\qquad 0\le e_i<t.
\]
On $x_0\ne0$, these monomials form a basis of the upper function field
over the field generated by the $t$-th powers of the affine coordinates.
A linear relation among their products with pulled-back sections
therefore separates into linear relations among the original sections,
which proves \eqref{eq:rank:power-sections}.

Apply \eqref{eq:rank:preliminary-count} to $\mathcal G_t$, whose highest
general-line degree is at most $t\mu+b(t-1)+u$.  Divide the resulting
bound by $t^b$ and let $t\to\infty$, keeping $\mathcal G$ and $u$ fixed.
This yields
\[
  h^0(\mathcal G)\le \frac{u}{b!}(\mu+b)^b,
\]
and hence the asserted bound, for example with $C_b=b^b/b!$.
\end{proof}

\begin{proof}[Completion of the proof of Theorem~\ref{thm:rank}]
Choose a finite linear projection $f:\widehat Y\to\PP^b$ for the
embedding defined by $H$.  It has degree $N=H^b$ and
$f^*\OO_{\PP^b}(1)=\OO_{\widehat Y}(H)$.
The torsion-free sheaf $f_*\mathcal E_m$ has rank $Nr_m$ and
\begin{equation}\label{eq:rank:pushforward-slope}
  \mu_{\max}(f_*\mathcal E_m)\le Cm\delta/N.
\end{equation}
For this, let $\mathcal A$ be its first Harder--Narasimhan term.
Adjunction gives a nonzero map $f^*\mathcal A\to\mathcal E_m$.
After removing torsion and extending across codimension two, it gives
a map from the reflexive pullback $f^{[*]}\mathcal A$ to the reflexive
sheaf $\mathcal E_m$; denote its torsion-free image by $\mathcal Q$.
Finite separable reflexive pullback preserves semistability and
multiplies slopes for these polarizations by $N$
\cite[Lemma~2.1]{Weissmann2024}.
Consequently
\[
  N\mu(\mathcal A)
   \le\mu_H(\mathcal Q)
   \le\mu_{\max,H}(\mathcal E_m)
   \le Cm\delta,
\]
proving \eqref{eq:rank:pushforward-slope}.  No injectivity of the
adjunction map is required.

Apply Lemma~\ref{lem:rank:section-count} to the semistable
Harder--Narasimhan factors of $f_*\mathcal E_m$ and sum their section
bounds.  Their ranks sum to $Nr_m$, and $N$ is bounded in terms of $n$.
For all sufficiently large divisible $m$, this gives
\[
  \dim V_m\le h^0(\widehat Y,\mathcal E_m)
       \le Cr_m(m\delta)^b
       \le C\bigl(m\eta v(D)\bigr)^s(m\delta)^b.
\]
This is the target estimate \eqref{eq:rank:count-target}.
Divide by $m^n$ and use the volume asymptotic in divisible degrees.
Since $n=s+b$, we obtain
\[
  \frac{v(D)^n}{n!}\le C\eta^s v(D)^s\delta^b.
\]
The lower calibration bound $v(D)\ge c\delta$ then implies
$c^b/n!\le C\eta^s$.  This is impossible for a sufficiently small
positive $\eta$ depending only on $n$.  There are only finitely many
pairs $b,s>0$ with $b+s=n$, so choose $\eta$ for all of them at once,
also satisfying the earlier verticality and chain requirements, and
then impose the rank threshold in Lemma~\ref{lem:rank:low-order}.
The resulting contradiction bounds $r(T)$ in terms of $n$ alone.
\end{proof}

